\documentclass[a4paper,12pt,twoside]{memoir}

% Castellano
\usepackage[spanish,es-tabla]{babel}
\selectlanguage{spanish}
\usepackage[utf8]{inputenc}
\usepackage[T1]{fontenc}
\usepackage{lmodern} % Scalable font
\usepackage{microtype}
\usepackage{placeins}

\RequirePackage{booktabs}
\RequirePackage[table]{xcolor}
\RequirePackage{xtab}
\RequirePackage{multirow}


\usepackage{minted}
\usepackage{algpseudocode}

% Links
\PassOptionsToPackage{hyphens}{url}\usepackage[colorlinks]{hyperref}
\hypersetup{
	allcolors = {red}
}

% Ecuaciones
\usepackage{amsmath}

% Rutas de fichero / paquete
\newcommand{\ruta}[1]{{\sffamily #1}}

% Párrafos
\nonzeroparskip

% Huérfanas y viudas
\widowpenalty100000
\clubpenalty100000

% Imagenes
\usepackage{graphicx}
\newcommand{\imagen}[2]{
	\begin{figure}[!h]
		\centering
		\includegraphics[width=0.9\textwidth]{#1}
		\caption{#2}\label{fig:#1}
	\end{figure}
	\FloatBarrier
}

\newcommand{\imagenflotante}[2]{
	\begin{figure}%[!h]
		\centering
		\includegraphics[width=0.9\textwidth]{#1}
		\caption{#2}\label{fig:#1}
	\end{figure}
}



% El comando \figura nos permite insertar figuras comodamente, y utilizando
% siempre el mismo formato. Los parametros son:
% 1 -> Porcentaje del ancho de página que ocupará la figura (de 0 a 1)
% 2 --> Fichero de la imagen
% 3 --> Texto a pie de imagen
% 4 --> Etiqueta (label) para referencias
% 5 --> Opciones que queramos pasarle al \includegraphics
% 6 --> Opciones de posicionamiento a pasarle a \begin{figure}
\newcommand{\figuraConPosicion}[6]{%
  \setlength{\anchoFloat}{#1\textwidth}%
  \addtolength{\anchoFloat}{-4\fboxsep}%
  \setlength{\anchoFigura}{\anchoFloat}%
  \begin{figure}[#6]
    \begin{center}%
      \Ovalbox{%
        \begin{minipage}{\anchoFloat}%
          \begin{center}%
            \includegraphics[width=\anchoFigura,#5]{#2}%
            \caption{#3}%
            \label{#4}%
          \end{center}%
        \end{minipage}
      }%
    \end{center}%
  \end{figure}%
}

%
% Comando para incluir imágenes en formato apaisado (sin marco).
\newcommand{\figuraApaisadaSinMarco}[5]{%
  \begin{figure}%
    \begin{center}%
    \includegraphics[angle=90,height=#1\textheight,#5]{#2}%
    \caption{#3}%
    \label{#4}%
    \end{center}%
  \end{figure}%
}
% Para las tablas
\newcommand{\otoprule}{\midrule [\heavyrulewidth]}
%
% Nuevo comando para tablas pequeñas (menos de una página).
\newcommand{\tablaSmall}[5]{%
 \begin{table}
  \begin{center}
   \rowcolors {2}{gray!35}{}
   \begin{tabular}{#2}
    \toprule
    #4
    \otoprule
    #5
    \bottomrule
   \end{tabular}
   \caption{#1}
   \label{tabla:#3}
  \end{center}
 \end{table}
}

%
% Nuevo comando para tablas pequeñas (menos de una página).
\newcommand{\tablaSmallSinColores}[5]{%
 \begin{table}[H]
  \begin{center}
   \begin{tabular}{#2}
    \toprule
    #4
    \otoprule
    #5
    \bottomrule
   \end{tabular}
   \caption{#1}
   \label{tabla:#3}
  \end{center}
 \end{table}
}

\newcommand{\tablaApaisadaSmall}[5]{%
\begin{landscape}
  \begin{table}
   \begin{center}
    \rowcolors {2}{gray!35}{}
    \begin{tabular}{#2}
     \toprule
     #4
     \otoprule
     #5
     \bottomrule
    \end{tabular}
    \caption{#1}
    \label{tabla:#3}
   \end{center}
  \end{table}
\end{landscape}
}

%
% Nuevo comando para tablas grandes con cabecera y filas alternas coloreadas en gris.
\newcommand{\tabla}[6]{%
  \begin{center}
    \tablefirsthead{
      \toprule
      #5
      \otoprule
    }
    \tablehead{
      \multicolumn{#3}{l}{\small\sl continúa desde la página anterior}\\
      \toprule
      #5
      \otoprule
    }
    \tabletail{
      \hline
      \multicolumn{#3}{r}{\small\sl continúa en la página siguiente}\\
    }
    \tablelasttail{
      \hline
    }
    \bottomcaption{#1}
    \rowcolors {2}{gray!35}{}
    \begin{xtabular}{#2}
      #6
      \bottomrule
    \end{xtabular}
    \label{tabla:#4}
  \end{center}
}

%
% Nuevo comando para tablas grandes con cabecera.
\newcommand{\tablaSinColores}[6]{%
  \begin{center}
    \tablefirsthead{
      \toprule
      #5
      \otoprule
    }
    \tablehead{
      \multicolumn{#3}{l}{\small\sl continúa desde la página anterior}\\
      \toprule
      #5
      \otoprule
    }
    \tabletail{
      \hline
      \multicolumn{#3}{r}{\small\sl continúa en la página siguiente}\\
    }
    \tablelasttail{
      \hline
    }
    \bottomcaption{#1}
    \begin{xtabular}{#2}
      #6
      \bottomrule
    \end{xtabular}
    \label{tabla:#4}
  \end{center}
}

%
% Nuevo comando para tablas grandes sin cabecera.
\newcommand{\tablaSinCabecera}[5]{%
  \begin{center}
    \tablefirsthead{
      \toprule
    }
    \tablehead{
      \multicolumn{#3}{l}{\small\sl continúa desde la página anterior}\\
      \hline
    }
    \tabletail{
      \hline
      \multicolumn{#3}{r}{\small\sl continúa en la página siguiente}\\
    }
    \tablelasttail{
      \hline
    }
    \bottomcaption{#1}
  \begin{xtabular}{#2}
    #5
   \bottomrule
  \end{xtabular}
  \label{tabla:#4}
  \end{center}
}



\definecolor{cgoLight}{HTML}{EEEEEE}
\definecolor{cgoExtralight}{HTML}{FFFFFF}

%
% Nuevo comando para tablas grandes sin cabecera.
\newcommand{\tablaSinCabeceraConBandas}[5]{%
  \begin{center}
    \tablefirsthead{
      \toprule
    }
    \tablehead{
      \multicolumn{#3}{l}{\small\sl continúa desde la página anterior}\\
      \hline
    }
    \tabletail{
      \hline
      \multicolumn{#3}{r}{\small\sl continúa en la página siguiente}\\
    }
    \tablelasttail{
      \hline
    }
    \bottomcaption{#1}
    \rowcolors[]{1}{cgoExtralight}{cgoLight}

  \begin{xtabular}{#2}
    #5
   \bottomrule
  \end{xtabular}
  \label{tabla:#4}
  \end{center}
}


















\graphicspath{ {./img/} }

% Capítulos
\chapterstyle{bianchi}
\newcommand{\capitulo}[2]{
	\setcounter{chapter}{#1}
	\setcounter{section}{0}
	\chapter*{#2}
	\addcontentsline{toc}{chapter}{#2}
	\markboth{#2}{#2}
}

% Apéndices
\renewcommand{\appendixname}{Apéndice}
\renewcommand*\cftappendixname{\appendixname}

\newcommand{\apendice}[1]{
	%\renewcommand{\thechapter}{A}
	\chapter{#1}
}

\renewcommand*\cftappendixname{\appendixname\ }

% Formato de portada
\makeatletter
\usepackage{xcolor}
\newcommand{\tutor}[1]{\def\@tutor{#1}}
\newcommand{\course}[1]{\def\@course{#1}}
\definecolor{cpardoBox}{HTML}{E6E6FF}
\def\maketitle{
  \null
  \thispagestyle{empty}
  % Cabecera ----------------
\noindent\includegraphics[width=\textwidth]{cabecera}\vspace{1cm}%
  \vfill
  % Título proyecto y escudo informática ----------------
  \colorbox{cpardoBox}{%
    \begin{minipage}{.8\textwidth}
      \vspace{.5cm}\Large
      \begin{center}
      \textbf{TFG del Grado en Ingeniería Informática}\vspace{.6cm}\\
      \textbf{\LARGE\@title{}}
      \end{center}
      \vspace{.2cm}
    \end{minipage}

  }%
  \hfill\begin{minipage}{.20\textwidth}
    \includegraphics[width=\textwidth]{escudoInfor}
  \end{minipage}
  \vfill
  % Datos de alumno, curso y tutores ------------------
  \begin{center}%
  {%
    \noindent\LARGE
    Presentado por \@author{}\\ 
    en Universidad de Burgos --- \@date{}\\
    Tutor: \@tutor{}\\
  }%
  \end{center}%
  \null
  \cleardoublepage
  }
\makeatother

\newcommand{\nombre}{Daniel Mellado Hurtado} %%% cambio de comando

% Datos de portada
\title{Monitorización sensores IOT}
\author{\nombre}
\tutor{Bruno Baruque Zanon}
\date{\today}

\begin{document}

\maketitle


\newpage\null\thispagestyle{empty}\newpage


%%%%%%%%%%%%%%%%%%%%%%%%%%%%%%%%%%%%%%%%%%%%%%%%%%%%%%%%%%%%%%%%%%%%%%%%%%%%%%%%%%%%%%%%
\thispagestyle{empty}


\noindent\includegraphics[width=\textwidth]{cabecera}\vspace{1cm}

\noindent D. nombre Bruno Baruque Zanon, profesor del departamento Ingeniería Informática, área de Ciencia de la computación e inteligencia.

\noindent Expone:

\noindent Que el alumno D. \nombre, con DNI 71363864Q, ha realizado el Trabajo final de Grado en Ingeniería Informática titulado Monitorización sensores IOT. 

\noindent Y que dicho trabajo ha sido realizado por el alumno bajo la dirección del que suscribe, en virtud de lo cual se autoriza su presentación y defensa.

\begin{center} %\large
En Burgos, {\large \today}
\end{center}

\vfill\vfill\vfill

% Author and supervisor
\begin{minipage}{0.45\textwidth}
\begin{flushleft} %\large
Vº. Bº. del Tutor:\\[2cm]
D. nombre tutor
\end{flushleft}
\end{minipage}
\hfill
\begin{minipage}{0.45\textwidth}
\begin{flushleft} %\large
Vº. Bº. del co-tutor:\\[2cm]
D. nombre co-tutor
\end{flushleft}
\end{minipage}
\hfill

\vfill

% para casos con solo un tutor comentar lo anterior
% y descomentar lo siguiente
%Vº. Bº. del Tutor:\\[2cm]
%D. nombre tutor


\newpage\null\thispagestyle{empty}\newpage




\frontmatter

% Abstract en castellano
\renewcommand*\abstractname{Resumen}
\begin{abstract}

El tratamiento, monitorización y gestión de datos procedentes de dispositivos inteligentes es más prioritario que nunca, poder reconocer patrones y tendencias en estos datos puede ayudar a predecir fallos o eventos no deseados con suficiente tiempo de antelación cómo para poder ser evitados. 

Con el fin de desarrollar un sistema capaz de mantener monitorizados y gestionados diversos dispositivos IOT se ha desarrollado esta aplicación.

\end{abstract}

\renewcommand*\abstractname{Descriptores}
\begin{abstract}
IoT, internet de las cosas, sensores, monitorización, inteligencia artificial, aprendizaje automático, aprendizaje incremental, predicciones, PRTG, Elasticsearch.
\end{abstract}

\clearpage

% Abstract en inglés
\renewcommand*\abstractname{Abstract}
\begin{abstract}

The processing, monitoring and management of data from smart devices is more of a priority than ever, being able to recognise patterns and trends in data can help to predict failures or unwanted events sufficiently in advance so that they can be avoided. 

In order to develop a system capable of keeping various IOT devices monitored and managed, this application has been developed.
\end{abstract}

\renewcommand*\abstractname{Keywords}
\begin{abstract}
IoT, internet of things, sensors, monitoring, artificial intelligence, machine learning, incremental learning, predictions, PRTG, Elasticsearch.
\end{abstract}

\clearpage

% Indices
\tableofcontents

\clearpage

\listoffigures

\clearpage

\listoftables
\clearpage

\mainmatter
\capitulo{1}{Introducción}


Vivimos en un mundo en el que todo está conectado y parametrizado, miles de dispositivos inteligentes capaces de captar y transmitir datos, desde los relojes que miden nuestras pulsaciones hasta los sensores que miden el nivel de oxígeno en la Estación Espacial Internacional.

Por ello el tratamiento de los datos procedentes de todos estos sensores es primordial hoy en día y poder reconocer patrones y tendencias en estos datos puede ser de gran utilidad, no solo para empresas sino para todo el mundo.

Con el fin de mantener monitorizados y gestionados diversos dispositivos IOT, los cuales han sido cedidos por un cliente interesado en el proyecto, se ha desarrollado esta aplicación.

A lo largo de esta memoria y sus anexos se verá con detalle cómo ha sido el desarrollo del proyecto, su diseñado y como utilizar el producto final. 

\section{Material adjunto}
Los materiales que se encuentran adjuntos en la entrega de este proyecto son los siguientes:
\begin{itemize}
    \item Máquina virtual ubuntu server con la aplicación ya instalada.
    \item Memoria del Trabajo de Fin de Grado más los documentos de anexos.
    \item Script de instalación \textit{install.sh}.
    \item Paquete debian \textit{Monitorizacion\_IOT\_1.0\_all.seb}.
    \item Fichero Jupyter Notebook con las pruebas de regresión realizadas.
    \item Ficheros con los datasets usados para las pruebas.
    \item El código del programa.
\end{itemize}

\capitulo{2}{Objetivos del proyecto}

El objetivo de este proyecto consta de la captación dinámica de diversos datos provenientes de sensores IOT, los datos, cedidos por el cliente se han de tratar y almacenar en una base de datos para poder ser consultados y monitorizados por el usuario en cualquier momento y a tiempo real. 

Se implementará también un algoritmo de aprendizaje automático que ha de aprender con los datos obtenidos a identificar patrones que puedan llevar a situaciones no deseadas y que esto pueda ayudar al usuario a tomar decisiones con la suficiente antelación.

\subsection{Objetivos}

\begin{enumerate}
    \item Utilizar un repositorio en github y aplicar conocimientos sobre gestión de proyectos.
    \item Realizar un algoritmo de aprendizaje automático para poder aplicar lo aprendido durante la carrera así cómo ampliar mis conocimientos sobre minería de datos.
    \item Adentrarme y aprender sobre tratamiento, procesado y gestión de datos mediante tecnologías como Stream Processing Frameworks.
    \item La implementación de motores de búsqueda para felicitar encontrar los datos deseados con la mayor precisión posible.
    \item Que la monitorización de los datos se muestre los más clara y accesible posible para el usuario.
\end{enumerate}

\capitulo{3}{Conceptos teóricos}

En aquellos proyectos que necesiten para su comprensión y desarrollo de unos conceptos teóricos de una determinada materia o de un determinado dominio de conocimiento, debe existir un apartado que sintetice dichos conceptos.

Algunos conceptos teóricos de \LaTeX \footnote{Créditos a los proyectos de Álvaro López Cantero: Configurador de Presupuestos y Roberto Izquierdo Amo: PLQuiz}.



\section{Secciones}

Las secciones se incluyen con el comando section.

\subsection{Subsecciones}

Además de secciones tenemos subsecciones.

\subsubsection{Subsubsecciones}

Y subsecciones. 


\section{Referencias}

Las referencias se incluyen en el texto usando cite \cite{wiki:latex}. Para citar webs, artículos o libros \cite{koza92}.


\section{Imágenes}

Se pueden incluir imágenes con los comandos standard de \LaTeX, pero esta plantilla dispone de comandos propios como por ejemplo el siguiente:

\imagen{escudoInfor}{Autómata para una expresión vacía}



\section{Listas de items}

Existen tres posibilidades:

\begin{itemize}
	\item primer item.
	\item segundo item.
\end{itemize}

\begin{enumerate}
	\item primer item.
	\item segundo item.
\end{enumerate}

\begin{description}
	\item[Primer item] más información sobre el primer item.
	\item[Segundo item] más información sobre el segundo item.
\end{description}
	
\begin{itemize}
\item 
\end{itemize}

\section{Tablas}

Igualmente se pueden usar los comandos específicos de \LaTeX o bien usar alguno de los comandos de la plantilla.

\tablaSmall{Herramientas y tecnologías utilizadas en cada parte del proyecto}{l c c c c}{herramientasportipodeuso}
{ \multicolumn{1}{l}{Herramientas} & App AngularJS & API REST & BD & Memoria \\}{ 
HTML5 & X & & &\\
CSS3 & X & & &\\
BOOTSTRAP & X & & &\\
JavaScript & X & & &\\
AngularJS & X & & &\\
Bower & X & & &\\
PHP & & X & &\\
Karma + Jasmine & X & & &\\
Slim framework & & X & &\\
Idiorm & & X & &\\
Composer & & X & &\\
JSON & X & X & &\\
PhpStorm & X & X & &\\
MySQL & & & X &\\
PhpMyAdmin & & & X &\\
Git + BitBucket & X & X & X & X\\
Mik\TeX{} & & & & X\\
\TeX{}Maker & & & & X\\
Astah & & & & X\\
Balsamiq Mockups & X & & &\\
VersionOne & X & X & X & X\\
} 

\capitulo{4}{Técnicas y herramientas}

Esta parte de la memoria tiene como objetivo presentar las técnicas metodológicas y las herramientas de desarrollo que se han utilizado para llevar a cabo el proyecto. Si se han estudiado diferentes alternativas de metodologías, herramientas, bibliotecas se puede hacer un resumen de los aspectos más destacados de cada alternativa, incluyendo comparativas entre las distintas opciones y una justificación de las elecciones realizadas. 
No se pretende que este apartado se convierta en un capítulo de un libro dedicado a cada una de las alternativas, sino comentar los aspectos más destacados de cada opción, con un repaso somero a los fundamentos esenciales y referencias bibliográficas para que el lector pueda ampliar su conocimiento sobre el tema.


\section{Stream Processing Framework}

Stream processing es una tecnología que procesa los datos de forma continua y secuencial, para realizar esta tarea se utilizan flujos de datos infinitos y sin límites de tiempo. Son extremadamente útiles a la hora de monitorizar sistemas, redes, aplicaciones, dispositivos, etc... 

algunas de las alternativas más populares que he considerado son las siguientes:


\subsection{Apache Kafka}

Apache Kafka es un sistema de intermediación de mensajes que responde al patrón "Publish/Subscribe Messaging" que se utiliza para la comunicación entre aplicaciones. Entre sus principales características se encuentran que es un sistema escalable y persistente, con gran tolerancia a fallos y gran velocidad tanto de escritura cómo de lectura.

Para transmitir estos datos Kafka crea Topics, flujos de datos que a su vez se dividen en particiones, cada mensaje se almacena en una de estas particiones.

\subsection{Apache Samza}
Es un procesador de flujo a tiempo real, usa Apache Kafka para la mensajería y Apache Hadoop YARN para tolerancia a fallos, seguridad, independencia de procesos y gestión de recursos para su utilización es necesario Hadoop. Samza utiliza streams inmutables e implementa el procesamiento de flujo sin y con estado.  

\subsection{Apache Flume}

Apache Flume es un servicio distribuido capaz de mover grandes cantidades de datos y logs, este forma parte del ecosistema de Hadoop. Su arquitectura es tanto flexible cómo sencilla aunque presenta cierta dificultad a la hora de escalar. Esta basado en flujo de datos en Streaming permitiendo múltiples flujos. 

\section{Visualización y análisis de datos}

\subsection{PRTG Network Monitor}

PRTG Network Monitor es una herramienta de monitorización de Paessler, nos permite supervisar y monitorizar datos de sensores,  la versión gratuita esta limitada a 100 sensores, una de las ventajas de PRTG es la cantidad de variedad de sensores que existen de forma predefinida. Nos permite mostrar alertas al encontrar problemas o métricas inusuales y estas nos pueden llegar de diversas formas.

\subsection{ThingSpeak}

ThinkSpeak es una plataforma open source orientada al internet de las cosas (IOT) que nos permite recoger, almacenar datos de sensores en la nube usando el protocolo HTTP y visualizar los datos. También nos permite realizar análisis de los datos usando MATLAB. 

ThingSpeak consta de aplicaciones complementarias con diversas funcionalidades.

\subsection{Grafana}

Grafana nos permite visualizar y representar métricas de datos sin importar donde estén almacenados, la visualización es rápida y flexible con multitud de opciones. En Grafana se pueden definir alertas sobre las métricas que deseemos y mandarnos notificaciones de las mismas, también cuenta con cientos de plugins oficiales. 

Contiene una opción de pago en la que ellos albergan ellos el servidor y otro completamente gratuita en la que nosotros podremos ejecutarlo en cualquiera de los sistemas operativos principales.

\section{Motor de búsquedas}

\subsection{Algolia}
Algolia es un motor de búsqueda accesible vía API, la cual se puede integrar en web y aplicaciones móviles. Funciona exclusivamente con datos en formato JSON. Algolia posee una inteligencia artificial que aprende del usuario para mejorar su experiencia de uso. Consta de una versión gratuita, la cual tiene limitada el número de resultados de búsqueda.

\subsection{ElectricSearch}
EleactricSearch es uno de los motor de búsqueda mas valorados del mercado, de código abierto que proporciona análisis y búsqueda de datos a tiempo real. Esta desarrollado en Java y sólo soporta JSON cómo tipos de respuesta.

Posee una arquitectura distribuida siendo la escalabilidad horizontal, una de sus ventajas es su rapidez a la hora de realizar búsquedas de textos complejos para facilitar esto funciona mediante índices invertidos.  



\subsection{Solr}
Apache Solr es un motor de búsqueda open source escrito en java con comandos escritos en HTTP y guarda los archivos utilizando XML. Esta optimizado para trafico de altos volúmenes de datos y todo en tiempo real. 

\section{librerias}

\subsection{Scikit-learn}

Scikit-learn es un módulo de programación en python para machine learning, nos permite realizar algoritmos de clasificación, regresión, Análisis de grupos, Máquinas de vectores de soporte, Árboles de decisión, etc...

Scikit-learn trabaja a su vez con otras librerías cómo NumPy, SciPy, pandas, SymPy, Matplotlib, Ip[y].

\subsection{scikit-multiflow}

Scikit-multiflow es una librería de machine learning para python dedicada a streaming data y multi-output learning que nos permite generar y evaluar data streams.

\capitulo{5}{Aspectos relevantes del desarrollo del proyecto}


En este apartado se detallarán los aspectos más relevantes del proyecto.

\section{Motivación en la elección del proyecto}

La motivación a la hora de elegir este proyecto surgió por lo interesante de la premisa, tratar con datos de sensores reales y diseñar un sistema que no solo mantenga monitorizados dichos sensores, sino que aprenda de ellos resulto ser suficiente para la elección de este proyecto. 

\section{Desarrollo del proyecto}

\subsection{Inicio}

Las primeras semanas tras la asignación del proyecto se destinaron en asentar la idea general del proyecto y a establecer una serie de objetivos, así como la ruta que había que llevar para lograrlos con éxito.

Durante esta primeara etapa, la mayoría del tiempo se destinó al aprendizaje y a la investigación de posibles herramientas y tecnologías que podrían ser de utilidad.

Se eligió Elasticsearch como principal motor de búsqueda y como base de datos para los sensores. A partir de ahí comenzó el desarrollo del proyecto.

\subsection{Reuniones con el cliente}

A lo largo del proyecto se celebraron una serie de reuniones con el cliente, el cual dispone de una serie de sensores encargados de medir ciertos parámetros de una bodega. Estos sensores los tenía conectados al software de monitorización PRTG, por lo tanto, se estableció como punto de partida y cómo uno de los primeros objetivos el recopilar los datos de PRTG. 

\subsection{Conexión entre servidores}
El primer gran reto que se presentó fue realizar la conexión entre el servidor donde el cliente tiene albergados los sensores, PRTG y el servidor Ubuntu , el cual contenía Elasticsearch que se encargaría de la monitorización de dichos sensores. 

A la hora de descargar los datos de PRTG surgieron una serie de adversidades que se fueron solucionando, estos inconvenientes se encuentran detallados en el apartado \ref{cap:Problemas}. Una vez se logró que se fueran descargando los datos procedentes de PRTG de forma periódica llego la hora de configurar las herramientas de \textit{ELK Stack} para poder tratar los datos obtenidos y poder ser indexados correctamente en Elasticsearch.

Una vez conseguida la conexión entre los dos sistemas y logrado el mantener un flujo continuo de datos también se logró uno de los objetivos propuestos, monitorizar los datos. Gracias a Elasticsearch y Kibana se podía mantener un control sobre los datos de los sensores pudiendo representarlo de diversas formas gráficas.

\subsection{Aplicación de modelos machine learning}

El siguiente gran reto fue la realización de un modelo capaz de poder aprender de los datos recibidos y así poder realizar predicciones a futuro.

Durante esta etapa se llevó a cabo la investigación sobre que métodos y algoritmos utilizar. Se llegó a la conclusión de que la mejor idea y la más interesante, sería implementar algoritmos de aprendizaje incremental, de esta forma el modelo que se cree podrá aprender según obtenga los datos. Para ello se decidió la utilización de la librería python River, la cual, pese a seguir en desarrollo y contar con una documentación aún por terminar, es una librería prometedora.

Se diseñó un sistema con el cual tras descargarse de los datos de PRTG y almacenarlos en Elasticsearch se volvían a extraer para poder ser entrenados por un modelo. Se decidió que se almacenaran y luego se volvieran a extraer en vez de utilizar los datos obtenidos directamente de PRTG porque de esta forma el entrenamiento no dependía de PRTG y si se desactiva esta funcionalidad durante un tiempo y se vuelve a activar después, poder tener la opción de entrenar el modelo con datos pasados.

\subsection{Facilitar la instalación}
Para terminar, debido a la compleja configuración de todos los sistemas involucrados se decidió realizar un instalador que facilitará al usuario poner en marcha el programa.

\section{Problemas y resoluciones}\label{cap:Problemas}
A continuación, se mostrarán una serie de problemas que se fueron presentando durante todo el desarrollo del proyecto, así como la solución a la que se llegó.

\begin{itemize}
    \item El primer inconveniente que se presento fue a la hora de descargar los datos de PRTG ya que no permite descargar datos en intervalos menores a media hora. Para solucionar este problema y poder descargar los datos de forma más continua, se decidió descargar desde PRTG en intervalos de media hora pero realizar esta acción con una periodicidad de un minuto, de esta forma la primera vez que se ejecuta el programa se suben los datos descargados en su totalidad pero en los ciclos posteriores el nuevo conjunto de datos descargados de PRTG (con un rango de media hora) es comparado con el del ciclo anterior eliminando los datos repetidos, de esta forma se logra sacar los datos correspondientes al último minuto.
    
    
    \item Al descargarse los datos desde PRTG estos estos se encontraban con un formato difícil de procesar por logstash (ya que este está pensado principalmente para trabajar con logs). Por ello se decidió realizar un sencillo programa en flex que transformara las líneas de datos en un formato legible para logstash.\ref{cap:TransformacionDatos}
    
    Tras unos meses, se vio necesario realizar un tratamiento del campo fecha, La horas descargadas por PRTG vienen en un formato H:mm:ss mientras que Logstash lee las horas con un formato HH:mm:ss por lo que al encontrarse horas de un solo dígito logstash no lo reconocía y saltaba una excepción.
    
    se utilizó flex en un primer lugar por su sencillez, pero tras esta nueva necesidad se vio necesario migrar esta funcionalidad a python y prescindir del programa realizado en flex.
    
    \item Elasticsearch almacena las variables de tiempo en formato UTC pero a la ahora de visualizar los datos en kibana se utiliza por defecto la zona horaria del navegador.
    
    Una forma de solucionarlo es configurar Kibana para que cambiar la zona horaria a UTC.
    
    \item PRTG solo recoge los datos de los sensores cuando este se encuentre en un equipo encendido y con acceso a internet, si no, los datos procedentes de los sensores se pierden, haciendo que nunca se guardarán datos de varios días seguidos.
    
    Para solucionar esto se debería migrar el sistema a un servidor que se mantenga funcional siempre.
\end{itemize}

\capitulo{6}{Trabajos relacionados}

Este apartado sería parecido a un estado del arte de una tesis o tesina. En un trabajo final grado no parece obligada su presencia, aunque se puede dejar a juicio del tutor el incluir un pequeño resumen comentado de los trabajos y proyectos ya realizados en el campo del proyecto en curso. 

\capitulo{7}{Conclusiones y Líneas de trabajo futuras}

Para concluir, en este último apartado se compartirán las conclusiones a las que se ha llegado tras la realización de este proyecto, así como unas líneas de trabajo futuras que servirán para mejorar el proyecto.

\section{Conclusión}

Tras la finalización del proyecto se entrega un producto totalmente funcional, capaz de monitorizar diversos sensores y lograr realizar predicciones a futuro, pese a que las predicciones no sean tan fiables cómo cabría de esperar, este es sin duda un apartado en el que habría que mejorar, modificar los algoritmos o buscar nuevas soluciones que pudieran realizar predicciones fiables.

Durante el transcurso del proyecto no solo se han logrado los objetivos generales sino los personales también, el realizar un proyecto cómo este para lo bueno y para lo malo ha sido una experiencia enriquecedora tanto profesionalmente como personalmente.


\section{Líneas de trabajo futuras}
Esta aplicación se podría mejorar con la implementación nuevas funcionalidades, así como la mejora de las ya existentes.

A continuación, se listan algunas de ellas:
\begin{itemize}
    \item Mover el sistema a un servidor real, de esta manera se podrá captar los datos de los sensores de forma continuada y sin pérdidas.
    \item Implementar una interfaz de usuario para facilitar la experiencia del usuario. Mediante dicha interfaz se propone la implementación de los siguientes apartados:
    \begin{itemize}
        \item Añadir y borrar sensores de manera gráfica.
        \item Poder crear modelos para los sensores de forma intuitiva.
    \end{itemize}
    \item Guardar metadatos de los sensores que se vayan a introducir en Elasticsearch como pueden ser el nombre del sensor, localización, etc. 
    \item Mejorar los algoritmos de aprendizaje incremental así cómo añadir nuevos modelos capaces de realizar mejores predicciones.
\end{itemize}


\bibliographystyle{plain}
\bibliography{bibliografia}

\end{document}
